\chapter{Constitutive Relations Recovered by OpenSG}
\label{ch:constitutive}

Every OpenSG run answers the same question. Given a structure gene (SG) --- a
one-dimensional through-thickness stack, a two-dimensional cross-section or
cell, or a three-dimensional unit cell --- what is the constitutive law of the
macroscopic model that replaces it? The SG problem is solved once, and its
output is a small dense stiffness matrix relating \emph{generalized forces} to
\emph{generalized strains}. This chapter lists every such law the code
produces, the exact ordering of the strain and force components in each, the
function that computes it, the input file you edit, the command you type, and
the file the answer is written to.

The map is short. Every law in this chapter is one path through
the tree below: an SG of a given dimension, a model choice made either by the
\texttt{n\_model} argument of the solid engine or by which
\texttt{opensg\_shell} entry point is called, the macro law that comes back, and
the \texttt{.out} file it is written to. Solid lines are the default path of
each entry point; dashed lines are the alternative model choices that reuse the
same input file.

\captree{
\node[funcnode, text width=41mm] (f0)
  {\texttt{opensg\_solid.rm\_plate\_1D}\\ \texttt{.msg\_rm\_plate.rm\_plate\_msg}\\
   \texttt{model:} 0 or 1};
\node[inputnode, text width=29mm, left=5mm of f0] (i0)
  {\texttt{layup\_db.yaml}\\ 1-D SG: one laminate};
\node[datanode, text width=25mm, right=5mm of f0] (d0)
  {ABD $6\times6$ or\\ ABDG $8\times8$};
\node[outnode, text width=33mm, right=5mm of d0] (o0)
  {\texttt{<db>\_plate\_homo.out}};
\node[funcnode, text width=41mm, below=of f0] (f1)
  {\texttt{opensg\_solid.sg\_homo}\\ \texttt{.plate\_homo\_2d},
   \texttt{n\_model=1}\\ $\to$ \texttt{\_beam\_homo\_kkt}
   (needs \texttt{n\_sg} $\ge$ 2)};
\node[funcnode, text width=41mm, below=of f1] (f2)
  {\texttt{opensg\_solid.sg\_homo}\\ \texttt{.plate\_homo\_2d},
   \texttt{n\_model=2}\\ optionally \texttt{shear\_refined=True} $\to$
   \texttt{plate\_shear\_ladder}};
\node[funcnode, text width=41mm, below=of f2] (f3)
  {\texttt{opensg\_solid.sg\_homo}\\ \texttt{.plate\_homo\_2d},
   \texttt{n\_model=3}};
\node[inputnode, text width=29mm, left=5mm of f2] (i1)
  {solid SG\\ \texttt{<sg>.sc} or \texttt{<sg>.yaml}\\ \texttt{n\_sg} = 1, 2 or 3};
\node[datanode, text width=25mm, right=5mm of f1] (d1)
  {Timoshenko $6\times6$\\ (+ EB $4\times4$)};
\node[datanode, text width=25mm, right=5mm of f2] (d2)
  {ABD $6\times6$ or\\ ABDG $8\times8$};
\node[datanode, text width=25mm, right=5mm of f3] (d3)
  {3-D solid $6\times6$\\ + eng.\ constants};
\node[outnode, text width=33mm, right=5mm of d1] (o1) {\texttt{<sg>.out}};
\node[outnode, text width=33mm, right=5mm of d2] (o2) {\texttt{<sg>.out}};
\node[outnode, text width=33mm, right=5mm of d3] (o3) {\texttt{<sg>.out}};
\node[funcnode, text width=41mm, below=of f3] (f4)
  {\texttt{opensg\_shell.sg\_homo}\\ \texttt{.build\_rm\_bundle} $\to$
   \texttt{ring\_indep}};
\node[funcnode, text width=41mm, below=of f4] (f5)
  {\texttt{opensg\_shell.sg\_homo}\\ \texttt{.build\_solid\_bundle} $\to$
   \texttt{ring\_solid}};
\node[inputnode, text width=29mm, left=5mm of f4] (i2)
  {1-D shell SG\\ \texttt{iea\_s10\_shell.yaml}\\ cross-section contour};
\node[datanode, text width=25mm, right=5mm of f4] (d4)
  {Timoshenko $6\times6$\\ + wall $8\times8$ per section};
\node[outnode, text width=33mm, right=5mm of d4] (o4)
  {\texttt{<yaml>\_Timo.out}\\ \texttt{<yaml>\_ABDG.out}};
\node[datanode, text width=25mm, right=5mm of f5] (d5) {3-D solid $6\times6$};
\node[outnode, text width=33mm, right=5mm of d5] (o5)
  {\texttt{<yaml>\_C3D.out}};
\node[funcnode, text width=41mm, below=of f5] (f6)
  {\texttt{opensg\_shell.sg\_homo}\\ \texttt{.shell\_sg3d}};
\node[funcnode, text width=41mm, below=of f6] (f7)
  {\texttt{opensg\_shell.sg\_homo}\\ \texttt{.segment\_timo\_from\_3dyaml}};
\node[inputnode, text width=29mm, left=5mm of f6] (i3)
  {3-D shell SG \texttt{.yaml}\\ surface mesh, or\\ tapered segment};
\node[datanode, text width=25mm, right=5mm of f6] (d6) {3-D solid $6\times6$};
\node[outnode, text width=33mm, right=5mm of d6] (o6)
  {\texttt{<yaml>\_C3D.out}};
\node[datanode, text width=25mm, right=5mm of f7] (d7)
  {Timoshenko $6\times6$};
\node[outnode, text width=33mm, right=5mm of d7] (o7)
  {\texttt{<yaml>\_Timo.out}};
\draw[capedge] (i0) -- (f0);  \draw[capedge] (f0) -- (d0);
\draw[capedge] (d0) -- (o0);
\draw[capedge2] (i1) -- (f1); \draw[capedge] (i1) -- (f2);
\draw[capedge2] (i1) -- (f3);
\draw[capedge] (f1) -- (d1);  \draw[capedge] (d1) -- (o1);
\draw[capedge] (f2) -- (d2);  \draw[capedge] (d2) -- (o2);
\draw[capedge] (f3) -- (d3);  \draw[capedge] (d3) -- (o3);
\draw[capedge] (i2) -- (f4);  \draw[capedge2] (i2) -- (f5);
\draw[capedge] (f4) -- (d4);  \draw[capedge] (d4) -- (o4);
\draw[capedge] (f5) -- (d5);  \draw[capedge] (d5) -- (o5);
\draw[capedge] (i3) -- (f6);  \draw[capedge2] (i3) -- (f7);
\draw[capedge] (f6) -- (d6);  \draw[capedge] (d6) -- (o6);
\draw[capedge] (f7) -- (d7);  \draw[capedge] (d7) -- (o7);
}

Every one of those files is printed by \texttt{write\_sc\_K}, except the
per-section wall law \texttt{<yaml>\_ABDG.out}, which
\texttt{write\_abdg\_out} writes block by block in the same number format. In
both cases the title above the matrix names the law, so the tree above and
Table~\ref{tab:outnames} together identify any output file the code produces.

Two conventions hold throughout and are worth internalising before reading any
matrix.

\begin{itemize}
\item \textbf{Units follow the input.} The single writer
  \texttt{write\_sc\_K} performs no unit conversion, so an SG given in metres
  with moduli in pascals returns N, N$\cdot$m and Pa, while an SG given in
  millimetres with moduli in MPa returns the corresponding mixed units (N/mm,
  N$\cdot$mm).
\item \textbf{Shears are engineering shears.} Every off-diagonal strain
  component carries the factor of two ($2\varepsilon_{12}$,
  $2\bar\Gamma_{23}$, $2\gamma_{13}$), which is why each shear stiffness
  appears once, not four times, in the quadratic strain energy.
\end{itemize}

\section{The \texttt{.out} contract}
\label{sec:out-contract}

Every homogenization in OpenSG-2.0 writes a timed text file in the SwiftComp
\texttt{.K} layout through one function,
\texttt{opensg\_solid.sg\_homo.write\_sc\_K}:

\begin{lstlisting}[style=opensg,caption={The single output writer. Every law in this chapter is printed by this call.}]
write_sc_K(path, C, solve_time=None, model="", constants=True, name="")
\end{lstlisting}

The file opens with an \verb| OpenSG <model>| banner, prints the effective
stiffness $\mathbf C$, then its exact inverse $\mathbf S = \mathbf C^{-1}$ (the
effective compliance), and closes with \verb| Time taken: <t> sec|. Numbers are
written by the helper \texttt{\_knum}, a 19-character right-justified
\verb|%.7E| mantissa with a signed three-digit exponent, which is exactly the
SwiftComp \texttt{.K} column format.

The \texttt{name} argument is the only thing that changes the matrix title, and
it identifies which macro law you are looking at (Table~\ref{tab:outnames}).

\begin{table}[htbp]
\centering
\caption{The \texttt{name} argument of \texttt{write\_sc\_K} and the title it
prints. The title is the fastest way to tell which law an \texttt{.out} file
holds.}
\label{tab:outnames}
\footnotesize
\begin{tabular}{lll}
\toprule
\texttt{name} passed to \texttt{write\_sc\_K} & title printed in the \texttt{.out} & macro law \\
\midrule
\texttt{"Timoshenko"} & \texttt{The Effective Timoshenko Stiffness Matrix} & beam $6\times6$ \\
\texttt{"Classical Plate"} & \texttt{The Effective Classical Plate Stiffness Matrix} & plate ABD $6\times6$ \\
\texttt{"Reissner-Mindlin Plate"} & \texttt{The Effective Reissner-Mindlin Plate Stiffness Matrix} & plate ABDG $8\times8$ \\
\texttt{""} (omitted) & \texttt{The Effective Stiffness Matrix} & 3-D solid $6\times6$ \\
\bottomrule
\end{tabular}
\end{table}

Only the last of these also prints
\verb| The Engineering Constants (Approximated as Orthotropic)|. That is the
\texttt{constants=True} branch, and it is meaningful only for a
three-dimensional elasticity law; the beam and plate calls pass
\texttt{constants=False}.

\section{Beam: the Timoshenko $6\times6$}
\label{sec:timo}

The beam macro model carries six generalized strains: one axial extension, two
transverse shears, one twist and two bending curvatures. In code order --- the
order stated in the module header of \texttt{src/opensg\_solid/sg\_homo.py} ---
these are

\begin{equation}
\bar{\boldsymbol\epsilon} \;=\;
\left\{\;\bar\epsilon_{11}\;\;\gamma_{12}\;\;\gamma_{13}\;\;
\kappa_{1}\;\;\kappa_{2}\;\;\kappa_{3}\;\right\}^{\mathsf T},
\end{equation}

and the conjugate generalized forces are the axial force, the two transverse
shear forces, the torque and the two bending moments --- the VABS ordering
$\{F_1\;F_2\;F_3\;M_1\;M_2\;M_3\}$ in which the shipped blade load table
\texttt{ff51\_rmc\_reform.dat} is written (its header reads
\verb|# eta F1 F2 F3 M1 M2 M3 (RM CENTER-ref BeamDyn, VABS order)|). The law is

\widematrix{
\begin{Bmatrix} F_1 \\ F_2 \\ F_3 \\ M_1 \\ M_2 \\ M_3 \end{Bmatrix}
=
\begin{bmatrix}
S_{11} & S_{12} & S_{13} & S_{14} & S_{15} & S_{16} \\
S_{12} & S_{22} & S_{23} & S_{24} & S_{25} & S_{26} \\
S_{13} & S_{23} & S_{33} & S_{34} & S_{35} & S_{36} \\
S_{14} & S_{24} & S_{34} & S_{44} & S_{45} & S_{46} \\
S_{15} & S_{25} & S_{35} & S_{45} & S_{55} & S_{56} \\
S_{16} & S_{26} & S_{36} & S_{46} & S_{56} & S_{66}
\end{bmatrix}
\begin{Bmatrix} \bar\epsilon_{11} \\ \gamma_{12} \\ \gamma_{13} \\
\kappa_{1} \\ \kappa_{2} \\ \kappa_{3} \end{Bmatrix}.
}

The matrix is fully populated in general --- a composite blade section couples
extension to bending and twist --- but its diagonal carries the familiar
engineering meaning

\begin{equation}
\mathrm{diag}(\mathbf S) = \left[\;EA\;\;GA_2\;\;GA_3\;\;GJ\;\;EI_2\;\;EI_3\;\right],
\label{eq:vabsdiag}
\end{equation}

which is how the drivers label their printout: the list
\verb|LBL = ["EA", "GA2", "GA3", "GJ", "EI2", "EI3"]| appears verbatim in
\texttt{examples/OpenSG\_shell/IEA\_blade\_beam/1\_homo\_beam\_props.py}. Read
Eq.~\eqref{eq:vabsdiag} as a naming convention, not as a claim of
decoupling. $S_{44}$ is the torsional stiffness only to the extent that the
twist channel is uncoupled; when $S_{45}$ or $S_{46}$ is large, $GJ$ as
conventionally understood is not $S_{44}$.

\subsection{Running it on a shell cross-section}

The msg-shell ring route homogenizes a cross-section contour with the
constrained six-DOF shell element (independent drilling $\omega_3$ enforced by
element-constant Lagrange multipliers, $\gamma_{23}$ tied by MITC) and takes
its wall transverse shear from the MSG least-squares plate solution. The user
edits one file, the one-dimensional shell SG YAML, which carries the contour
nodes, the line elements, the per-element orientations, the layup sections, the
materials and the \texttt{reference:} field. From
\texttt{examples/OpenSG\_shell/IEA\_blade\_beam}:

\begin{lstlisting}[style=shellst]
python 1_homo_beam_props.py
\end{lstlisting}

The driver's user block is three lines long; only the station name changes:

\begin{lstlisting}[style=opensg,caption={The complete user input of the shell beam driver. build\_rm\_bundle reads iea\_s10\_shell.yaml.}]
name = "iea_s10"                  # IEA-22 station at r/R = 0.20
B = build_rm_bundle(name + "_shell.yaml")
C6 = np.asarray(B["Timo"])
\end{lstlisting}

\texttt{build\_rm\_bundle} writes two files by itself: the SwiftComp-layout
\texttt{<yaml base>\_Timo.out} (here \texttt{iea\_s10\_shell\_Timo.out}) with
banner \verb| OpenSG msg-shell beam model [ext sh2 sh3 twist bend2 bend3]| and
title \texttt{The Effective Timoshenko Stiffness Matrix}, and the per-section
wall law \texttt{<yaml base>\_ABDG.out} described in
Section~\ref{sec:wallabdg}. The driver additionally saves the bare $6\times6$
as \texttt{iea\_s10\_beam\_Timo.out} and the ring mesh figure
\texttt{iea\_s10\_ring\_mesh.png}.

\begin{figure}[htbp]
\centering
\includegraphics[width=0.86\linewidth]{\FIGDIR/shell_xsec_mesh.png}
\caption{The one-dimensional shell SG solved by \texttt{build\_rm\_bundle}: the
IEA-22 blade cross-section at $r/R = 0.20$, 307 contour nodes and 310 two-node
line elements along the wall mid-surface, coloured by \texttt{sets: element}
entry. The Timoshenko $6\times6$ beam law whose diagonal is quoted in this
section is this mesh's answer, written to \texttt{iea\_s10\_shell\_Timo.out}.}
\label{fig:const-shellxsec}
\end{figure}

For the IEA-22 blade station at $r/R = 0.20$ --- the SG of
Figure~\ref{fig:const-shellxsec} --- that $6\times6$ has the diagonal

\widematrix{
\mathrm{diag}(\mathbf S) =
\left[\,2.7661\times10^{10},\;7.1861\times10^{8},\;4.2174\times10^{8},\;
2.4376\times10^{9},\;3.5144\times10^{10},\;6.8130\times10^{10}\,\right]
}

in SI units, so $EA = 27.7$ GN, $EI_3 = 68.1$ GN$\cdot$m$^2$ edgewise against
$EI_2 = 35.1$ GN$\cdot$m$^2$ flapwise, and $GJ = 2.44$ GN$\cdot$m$^2$. The
largest off-diagonal terms in that file are $S_{15} = 1.66864\times10^{9}$ and
$S_{16} = -2.07730\times10^{9}$ (extension--bending, i.e.\ the section
reference axis is not the tension centre) and
$S_{56} = -5.55712\times10^{9}$ (the two bending channels are not principal).
Terms that should vanish by symmetry sit at $10^{-5}$ against diagonals of
$10^{10}$, which is the numerical zero of this solve.

\subsection{Running it on a solid cross-section}

The msg-solid route runs the KKT beam driver: a zeroth-order $V_0$ solve under
four rigid-body Lagrange constraints, then a first-order $V_{1s}$ solve reusing
the same factorization, then the $6\times6$ reduction in
\texttt{finalize\_v1\_and\_compute\_deff}. The entry point is
\texttt{plate\_homo\_2d(..., n\_model=1)}, which requires \texttt{n\_sg >= 2}
(a two- or three-dimensional solid SG). From
\texttt{examples/OpenSG-solid/7\_get\_beam\_props\_from\_SG}:

\begin{lstlisting}[style=shellst]
python beam_homo_sg.py
\end{lstlisting}

\begin{lstlisting}[style=opensg,caption={The user block of beam\_homo\_sg.py. The material rows are (E1,E2,E3,G12,G13,G23,nu12,nu13,nu23) in MPa; angles rotates each material about the SG normal.}]
n_model = 1                    # 1: Beam; 2: Plate; 3: 3D elastic
name = "RHC_SW_2UC_45"         # reads <name>.yaml

material_param = jnp.array([
    (108e3, 8e3, 8e3, 4e3, 4e3, 3e3, 0.32, 0.32, 0.30),      # ply +45
    (108e3, 8e3, 8e3, 4e3, 4e3, 3e3, 0.32, 0.32, 0.30),      # ply -45
    (69e3, 69e3, 69e3, 26.54e3, 26.54e3, 26.54e3, 0.30, 0.30, 0.30)])
angles = jnp.array([45.0, -45.0, 0.0])

r = plate_homo_2d(name + ".yaml", material_param=material_param,
                  angles=angles, n_model=n_model)
\end{lstlisting}

This writes \texttt{RHC\_SW\_2UC\_45.out} with banner
\verb| OpenSG msg-solid beam model, omega 1| and the mesh figure
\texttt{RHC\_SW\_2UC\_45\_mesh.png} (Figure~\ref{fig:rhcmesh}). Its diagonal is

\widematrix{
\mathrm{diag}(\mathbf S) =
\left[\,9.8993\times10^{6},\;1.9075\times10^{6},\;9.3619\times10^{5},\;
2.9775\times10^{8},\;2.1896\times10^{9},\;9.2899\times10^{8}\,\right]
}

in the input units (N and N$\cdot$mm$^2$, since the SG is in millimetres and the
moduli in MPa).

\begin{figure}[htbp]
\centering
\includegraphics[width=0.72\linewidth]{\FIGDIR/rhc_2dsg_mesh.png}
\caption{The two-dimensional honeycomb solid SG \texttt{RHC\_SW\_2UC\_45}, run
here with $n_{\mathrm{model}} = 1$: the Timoshenko $6\times6$ beam law whose
diagonal is quoted in this section is this mesh's answer, written to
\texttt{RHC\_SW\_2UC\_45.out}. The same mesh is the SG of the plate
($n_{\mathrm{model}} = 2$) and solid ($3$) examples. Colours are the three
material sets of the \texttt{.yaml}: the $+45$ and $-45$ composite plies of the
cell walls and the isotropic face sheets.}
\label{fig:rhcmesh}
\end{figure}

\subsection{The classical $4\times4$ that rides along}
\label{sec:ceffeb}

The Timoshenko reduction is built on top of a zeroth-order (Euler--Bernoulli)
solve, and that intermediate law is a genuine classical beam stiffness in its
own right. It drops the two transverse shears:

\begin{equation}
\begin{Bmatrix} F_1 \\ M_1 \\ M_2 \\ M_3 \end{Bmatrix}
=
\mathbf A_{\mathrm{EB}}
\begin{Bmatrix} \bar\epsilon_{11} \\ \kappa_1 \\ \kappa_2 \\ \kappa_3 \end{Bmatrix},
\qquad \mathbf A_{\mathrm{EB}} \in \mathbb R^{4\times4}.
\end{equation}

On the solid route it is returned as \texttt{r["C\_eff\_EB"]} alongside the
Timoshenko \texttt{r["C\_eff"]}. In \texttt{\_beam\_homo\_kkt} it is formed as
$\mathbf C_{\mathrm{eb}} = (D_{ee} + D_1^{V_0})/\omega$ immediately after the
$V_0$ solve, before the first-order ladder runs, and it is then passed straight
into \texttt{finalize\_v1\_and\_compute\_deff} as the $\mathbf A$ block of the
Timoshenko transformation. It is not written to the \texttt{.out} --- the
\texttt{.out} always carries the $6\times6$ --- so a driver that wants it must
print it, as \texttt{beam\_homo\_sg.py} does.

\subsection{Which routes produce the beam law}

\begin{table}[htbp]
\centering
\caption{The three entry points that return a Timoshenko $6\times6$. The
msg-solid entry point lives in \texttt{opensg\_solid.sg\_homo}, the two
msg-shell ones in \texttt{opensg\_shell}.}
\label{tab:beamroutes}
\footnotesize
\begin{tabular}{lll}
\toprule
route & entry point & SG input \\
\midrule
msg-solid beam & \texttt{plate\_homo\_2d(..., n\_model=1)} & 2-D or 3-D solid SG (\texttt{n\_sg >= 2}) \\
msg-shell ring & \texttt{build\_rm\_bundle} (calls \texttt{ring\_indep}) & 1-D shell SG (cross-section contour) \\
msg-shell segment & \texttt{segment\_timo\_from\_3dyaml} & 3-D shell segment YAML (tapered or prismatic) \\
\bottomrule
\end{tabular}
\end{table}

The segment route replaces axial periodicity by Dirichlet data. Each end
cross-section is extracted topologically by \texttt{sg\_mesh.extract}, solved on
its own with \texttt{ring\_indep} to give that ring's $\mathbf C_6$, $V_0$ and
$V_1$, and those warping fields are then imposed on the segment's boundary
nodes; the segment energy is divided by the segment length
$L_z$. It writes \texttt{<yaml base>\_Timo.out} with the same
\texttt{Timoshenko} title and the banner
\verb|msg-shell tapered/aperiodic segment (boundary V0/V1 Dirichlet, L=<Lz>)|.
For a prismatic segment the result must reproduce the boundary ring's own
$6\times6$; that identity is the standing verification of the route.

\section{Plate: classical ABD $6\times6$ and Reissner--Mindlin ABDG $8\times8$}
\label{sec:plate}

\subsection{Classical (Kirchhoff--Love) plate law}

The classical plate model pairs three membrane stress resultants and three
moment resultants with three membrane strains and three curvatures:

\widematrix{
\begin{Bmatrix} N_{11} \\ N_{22} \\ N_{12} \\ M_{11} \\ M_{22} \\ M_{12} \end{Bmatrix}
=
\begin{bmatrix} \mathbf A & \mathbf B \\[2pt] \mathbf B^{\mathsf T} & \mathbf D \end{bmatrix}
\begin{Bmatrix} \varepsilon_{11} \\ \varepsilon_{22} \\ 2\varepsilon_{12} \\
\kappa_{11} \\ \kappa_{22} \\ \kappa_{12} \end{Bmatrix},
\qquad
\mathbf A, \mathbf B, \mathbf D \in \mathbb R^{3\times3}.
}

The row order \verb|[N11 N22 N12 M11 M22 M12]| is stated in the module header
of \texttt{src/opensg\_solid/sg\_homo.py} and echoed by the driver in
\texttt{examples/OpenSG-solid/3\_get\_plate\_props\_from\_2D\_SG}. The
$\mathbf B$ block is the extension--bending coupling; it vanishes only for a
laminate that is symmetric about the chosen reference surface, so it is a
property of the SG \emph{and} of where you put the reference. Run it with

\begin{lstlisting}[style=shellst]
python plate_homo_2dsg.py
\end{lstlisting}

which calls \texttt{plate\_homo\_2d(name + ".yaml", material\_param=...,
angles=..., n\_model=2)} and writes \texttt{RHC\_SW\_2UC\_45.out} with banner
\verb| OpenSG msg-solid plate model, omega 35.248001, periodic| and title
\texttt{The Effective Classical Plate Stiffness Matrix}. Here $\omega =
35.248001$ is the in-plane SG measure (the cell width), and the leading entries
are $A_{11} = 3.3586\times10^{5}$, $A_{22} = 1.1291\times10^{5}$,
$A_{66} = 1.0552\times10^{5}$ N/mm and $D_{11} = 9.0368\times10^{7}$,
$D_{22} = 5.7011\times10^{7}$, $D_{66} = 4.6026\times10^{7}$ N$\cdot$mm. The
$\mathbf B$ block is at the $10^{0}$ level against an $\mathbf A$ of $10^{5}$,
i.e.\ the cell is symmetric about its own mid-surface to solver precision.

\begin{figure}[htbp]
\centering
\includegraphics[width=0.72\linewidth]{\FIGDIR/rhc_2dsg_mesh.png}
\caption{The same two-dimensional honeycomb cell SG \texttt{RHC\_SW\_2UC\_45} as
Figure~\ref{fig:rhcmesh}, run here with $n_{\mathrm{model}} = 2$: the classical
plate ABD $6\times6$ whose leading entries are quoted in this section is this
mesh's answer, written to \texttt{RHC\_SW\_2UC\_45.out} under the title
\texttt{The Effective Classical Plate Stiffness Matrix}. Only the model choice
changes between this run and the beam run of
Section~\ref{sec:timo}; the SG is byte-identical.}
\label{fig:rhcplate}
\end{figure}

\subsection{Reissner--Mindlin plate law}

The shear-refined model adds the transverse-shear pair. The strain vector grows
to eight components and the resultant vector gains the two transverse shear
forces $Q_1$ and $Q_2$:

\widematrix{
\begin{Bmatrix} \mathbf N \\ \mathbf M \\ \mathbf Q \end{Bmatrix}
=
\begin{bmatrix}
\mathbf A & \mathbf B & \mathbf 0 \\[2pt]
\mathbf B^{\mathsf T} & \mathbf D & \mathbf 0 \\[2pt]
\mathbf 0 & \mathbf 0 & \mathbf G
\end{bmatrix}
\begin{Bmatrix} \boldsymbol\varepsilon \\ \boldsymbol\kappa \\ \boldsymbol\gamma \end{Bmatrix},
\qquad
\boldsymbol\gamma = \{\,2\gamma_{13}\;\;2\gamma_{23}\,\}^{\mathsf T},
\qquad
\mathbf Q = \{\,Q_1\;\;Q_2\,\}^{\mathsf T}.
}

Here $\boldsymbol\varepsilon$ and $\boldsymbol\kappa$ are exactly the first and
second triples of the classical law above.

The zero off-diagonal blocks are structural, not an approximation of
convenience. The code assembles the $8\times8$ as \verb|ABDG[:6,:6] = A6| and
\verb|ABDG[6:,6:] = G_msg|, so the in-plane law and the transverse-shear law
come from two separate variational problems and are then stacked. The
$2\times2$ block $\mathbf G$ is obtained as $\mathbf G = \mathbf X^{-1}$ from
the $\mathbf U^*$ least-squares reduction of the first-order warping ladder
(function \texttt{\_rm\_ls\_reduction}), a 78-equation, 27-unknown system solved
by a column-equilibrated truncated-SVD minimum-norm solve with cutoff
\texttt{\_LS\_RCOND = 1e-12}. Two diagnostics come out with it: the relative
residual of the fit, reported as \texttt{r["Ustar\_rel"]}, and the smallest
eigenvalue of $\mathbf X$. If that eigenvalue is not positive the SPD gate
fails and the block is returned as \texttt{None} --- in which case no
$8\times8$ exists and the code falls back to printing the classical
$6\times6$. Always check that you got an $8\times8$ before trusting a
transverse-shear number.

\subsubsection{The 1-D SG route}

\texttt{opensg\_solid.rm\_plate\_1D.msg\_rm\_plate.rm\_plate\_msg} solves the
through-thickness SG of a single laminate and returns \texttt{A6} (the
$6\times6$ ABD), \texttt{G\_msg} and \texttt{ABDG}, with rows
\verb|e11 e22 g12 k11 k22 k12 2g13 2g23|. Which of the two matrices is printed
is selected by \texttt{model:} in the user's \texttt{layup\_db.yaml}:
\texttt{model: 0} prints the classical $6\times6$, \texttt{model: 1} prints the
full $8\times8$. From
\texttt{examples/OpenSG-solid/1\_get\_plate\_props\_from\_1DSG}:

\begin{lstlisting}[style=shellst]
python 1d_sg.py
\end{lstlisting}

The input the user edits is \texttt{layup\_db.yaml} alone. It carries the
reference surface (\texttt{fraction: 0.5} is the mid-surface, \texttt{0} the
bottom or OML face, \texttt{1} the top), the materials with their densities,
the stacking sequence bottom-ply-first, and the SG discretisation
(\texttt{n\_per\_layer: 1} element per physical ply, \texttt{elem\_order: 4} for
the five-noded quartic element that represents the warping ladder exactly).
The script writes the SG mesh \texttt{1dsg.yaml} plus
\texttt{1dsg.png} (Figure~\ref{fig:plate1dsg}) and the law in
\texttt{layup\_db\_plate\_homo.out}.

The shipped case is the Nayak Example-5 sandwich, $(0/90/0/90/\mathrm{core})_s$,
nine plies, total thickness $h = 0.1524$ m, graphite/epoxy faces
($E_1 = 128$ GPa) over a HEREX C70.130 PVC foam core, each face ply
$0.9525$ mm and the core $144.78$ mm. Its \texttt{.out} banner records
\verb|9 plies, h = 0.152400 m, fraction = 0.5, rho*h = 30.251400 kg/m^2| and
the leading entries are

\begin{equation}
A_{11} = A_{22} = 5.49165\times10^{8}\ \mathrm{N/m},\quad
D_{11} = 3.00055\times10^{6},\quad D_{22} = 2.93711\times10^{6}\ \mathrm{N\!\cdot\!m},
\end{equation}
\begin{equation}
G_{11} = 7.70098\times10^{6}\ \mathrm{N/m},\qquad
G_{22} = 7.54237\times10^{6}\ \mathrm{N/m}.
\end{equation}

Read them physically: the thin stiff faces carry the membrane and bending
channels, while the thick compliant core sets the transverse shear. The
$\mathbf B$ block prints at $10^{-6}$ against an $\mathbf A$ of $10^{8}$,
because the stack is symmetric about \texttt{fraction: 0.5}; change that to
\texttt{fraction: 0.0} and $\mathbf B$ becomes a leading-order term.

\begin{figure}[htbp]
\centering
\includegraphics[width=0.62\textwidth]{\FIGDIR/plate_1dsg.png}
\caption{The through-thickness one-dimensional structure gene of the Nayak
Example-5 sandwich written by \texttt{1d\_sg.py}: eight graphite/epoxy face
plies over a HEREX C70.130 foam core, with the mid-surface reference
(\texttt{fraction: 0.5}) marked. This is the SG whose solution is the
Reissner--Mindlin ABDG $8\times8$ whose $A$, $D$ and $G$ entries are quoted in
this section, written by \texttt{rm\_plate\_msg} to
\texttt{layup\_db\_plate\_homo.out}.}
\label{fig:plate1dsg}
\end{figure}

\subsubsection{The 2-D / 3-D SG route}

\texttt{plate\_homo\_2d(..., n\_model=2, shear\_refined=True)} runs
\texttt{plate\_shear\_ladder} on the general SG assembly after the classical
solve, re-integrating the value-product ladder blocks on a quadrature exact to
degree $2p+2$. It adds \texttt{r["G\_msg"]} ($2\times2$),
\texttt{r["A6\_ladder"]}, \texttt{r["X\_shear"]}, \texttt{r["Ustar\_rel"]} and
\texttt{r["ABDG"]} to the result dict, and the \texttt{.out} switches its title
to \texttt{Reissner-Mindlin Plate} and its matrix to the $8\times8$. Passing
\texttt{shear\_refined=True} with any \texttt{n\_model} other than $2$ raises a
\texttt{ValueError}.

Note the derivation status recorded in the docstring of
\texttt{plate\_shear\_ladder}, because it decides how much to trust the
$\mathbf G$ block. For \texttt{n\_sg = 1} this \emph{is} the Yu (2005)
one-dimensional formulation, gated against the closed-form anchors (an
isotropic layer with $\nu = 0$ returns $\tfrac{5}{6}Gh$). For \texttt{n\_sg = 2}
or \texttt{3} it is the SwiftComp-style extension --- the same energy functional
with $\Gamma_h$ carrying the resolved-dimension derivatives, periodic
master--slave coupling, and the $\langle w \rangle = 0$ average constraint
replacing the one-dimensional node constraint. That extension is gated on the
laminate-as-strip equivalence (a laminate meshed as a 2-D or 3-D SG must
reproduce its 1-D $\mathbf G$ under refinement) but is \textbf{not} yet
validated against an external reference for SGs that are heterogeneous
\emph{in plane}. The shear-refined \emph{recovery} is likewise not ported:
plate dehomogenization stays the classical Kirchhoff-mode recovery.

\section{Solid: the equivalent 3-D law}
\label{sec:solid}

When the macro model is ordinary three-dimensional elasticity, the SG returns a
$6\times6$ anisotropic stiffness in the Voigt order $[11\;22\;33\;23\;13\;12]$
with engineering shear strains:

\widematrix{
\begin{Bmatrix}
\sigma_{11} \\ \sigma_{22} \\ \sigma_{33} \\ \sigma_{23} \\ \sigma_{13} \\ \sigma_{12}
\end{Bmatrix}
=
\mathbf C^{\,\mathrm{3D}}
\begin{Bmatrix}
\bar\Gamma_{11} \\ \bar\Gamma_{22} \\ \bar\Gamma_{33} \\
2\bar\Gamma_{23} \\ 2\bar\Gamma_{13} \\ 2\bar\Gamma_{12}
\end{Bmatrix},
\qquad \mathbf C^{\,\mathrm{3D}} \in \mathbb R^{6\times6}.
}

Internally the solid engine stores strain and stress in the SwiftComp order
$(xx, yy, zz, yz, xz, xy)$, which is the same list. On the msg-shell routes
axis 1 is the prismatic (beam) direction and axes 2 and 3 span the
cross-section, as recorded in the module constant
\texttt{opensg\_shell.sg\_homo.GBAR\_ORDER}:
\verb|strains [e11 e22 e33 2e23 2e13 2e12] -> stiffness C11..C66|
\verb|(1=beam axis, 2/3=cross-section axes)|.

\subsection{The engineering-constants block}

This is the only law for which the \texttt{.out} also prints an
engineering-constants block. It is computed from the compliance
$\mathbf S = (\mathbf C^{\,\mathrm{3D}})^{-1}$ by the orthotropic-approximation
formulas coded in \texttt{write\_sc\_K} (and mirrored on the shell side by
\texttt{opensg\_shell.sg\_homo.elastic\_constants}, which also returns
$\mathrm{cond}(\mathbf C)$):

\widematrix{
E_1 = \frac{1}{S_{11}},\quad E_2 = \frac{1}{S_{22}},\quad E_3 = \frac{1}{S_{33}},
\qquad
G_{12} = \frac{1}{S_{66}},\quad G_{13} = \frac{1}{S_{55}},\quad G_{23} = \frac{1}{S_{44}},
}
\widematrix{
\nu_{12} = -\frac{S_{12}}{S_{11}},\qquad
\nu_{13} = -\frac{S_{13}}{S_{11}},\qquad
\nu_{23} = -\frac{S_{23}}{S_{22}}.
}

The word \emph{approximated} in the printed heading is load-bearing. These nine
numbers reproduce $\mathbf C^{\,\mathrm{3D}}$ exactly only if that matrix is
genuinely orthotropic in the SG axes. A cell with normal--shear coupling still
prints them, and they will not rebuild the full matrix. The matrix, not the
constants block, is the result.

\subsection{Which entry points produce it}

\begin{table}[htbp]
\centering
\caption{The three entry points that return an equivalent 3-D solid law, and
the measure each divides the assembled energy by.}
\label{tab:solidroutes}
\footnotesize
\begin{tabular}{lll}
\toprule
entry point & SG & normalisation \\
\midrule
\texttt{plate\_homo\_2d(..., n\_model=3)} & 1-D/2-D/3-D solid SG & $\mathbf C_{\mathrm{eff}} = (\bar{\mathbf D} + \mathbf D_1)/\omega$ \\
\texttt{opensg\_shell.build\_solid\_bundle} & 1-D shell SG (cross-section contour) & $\mathbf C^{\,\mathrm{3D}} = \mathbf D_{\mathrm{eff}}/A_{\mathrm{cell}}$ \\
\texttt{opensg\_shell.sg\_homo.shell\_sg3d} & 3-D shell SG (surface mesh) & $\mathbf D_{\mathrm{eff}}/\omega$, $\omega$ = midsurface area \\
\bottomrule
\end{tabular}
\end{table}

For \texttt{build\_solid\_bundle} the normalising area defaults to the convex
hull of the contour (\texttt{area\_source = "hull"}); pass \texttt{cell\_area}
explicitly to use a different measure, and the banner records which was used.
The isotropic square-tube example in the folder
\texttt{examples/OpenSG\_shell/3\_get\_solid\_props\_from\_shell\_cross\_section/}
\texttt{square\_section\_isotropic}
sets \texttt{CELL\_AREA = 4*a*t\_w}, the wall material area, so that the shell
route and the solid route are normalised identically and can be compared term
by term. Its \texttt{square\_tube\_1Dshell\_C3D.out} returns
$E_1 = 7.0000000\times10^{10}$ Pa and $\nu_{12} = 0.30000$, recovering the wall
material's own $E$ and $\nu$ along the beam axis exactly, which is the sanity
check that fixes the normalisation.

The last row of Table~\ref{tab:solidroutes} carries a subtlety worth stating
plainly. \texttt{shell\_sg3d} \emph{returns} \texttt{C3D} divided by the SG
measure (the midsurface surface area, the three-dimensional analogue of the
plane-section perimeter), but \emph{writes} the \texttt{.out} divided by the
bounding-box cell volume $V_{\mathrm{cell}}$, so the printed moduli compare
directly against a solid SwiftComp \texttt{.K} file. If you read the numbers
out of the return dict and out of the \texttt{.out} and they differ by a
constant factor, that factor is $\omega/V_{\mathrm{cell}}$ and nothing is
wrong.

\subsection{Periodicity, and the aperiodic alternative}
\label{sec:periodicity}

A free SG driven by six macro strains is rank one by construction: every mode
except $\bar\Gamma_{11}$ can be cancelled at zero energy by an affine
fluctuation. Something must forbid those fields. All three routes
\emph{default} to periodicity, and periodicity is the SwiftComp-parity mode:

\begin{itemize}
\item \texttt{build\_solid\_bundle(..., periodic=True)} is the default; it ties
  opposite bounding-box faces through the shell periodic assembly map, and the
  repeated \texttt{dof\_map[dof\_map]} resolves the edge and corner chains, so
  faces, edges and corners are all tied. The tie rides in the assembly map, so
  no constraint rows are added and the rigid kernel drops from four rows to the
  three translations.
\item \texttt{shell\_sg3d(..., boundary="periodic")} is the default; the same
  map is applied on the full three-dimensional coordinates, and the three rigid
  translations are removed by area-weighted Lagrange rows.
\item \texttt{plate\_homo\_2d(..., boundary=None)} resolves to
  \texttt{"periodic"}.
\end{itemize}

Two of these also accept an aperiodic treatment, and it is a genuine
alternative, not a diagnostic. It is the \emph{boundary-solution} treatment:
for a unit cell the boundary solution is the macro (affine) field itself, and
mapping it onto the boundary prescribes zero fluctuation there.

\begin{itemize}
\item \texttt{shell\_sg3d(..., boundary="aperiodic")} clamps
  $w_1 = w_2 = w_3 = 0$ on every node of the bounding-box faces and leaves the
  three rotational DOFs \textbf{natural}. Translations are what the macro solid
  strain prescribes; clamping the rotations as well over-stiffens the cut
  edges, and the in-code note records the measured effect of freeing them
  ($+12\,\%$ down to $+7\,\%$ on $C_{11}$). No Lagrange border is needed,
  because the Dirichlet data also fixes the rigid modes.
\item \texttt{plate\_homo\_2d(..., boundary="aperiodic")} pins all three
  fluctuation DOFs of every bounding-box-face node --- the solid engine carries
  no rotational DOFs, so those three are the whole node --- instead of the usual
  first-node pin, and assembles on the raw mesh rather than the periodic map.
  It is supported for the plate and solid models with
  \texttt{solver="direct"}; requesting it with \texttt{n\_model=1} or
  \texttt{solver="cg"} raises a \texttt{ValueError}.
\end{itemize}

Kinematic (Dirichlet) homogenization is an \emph{upper bound} on a single cell,
and it converges to the periodic answer as the SG grows, because the boundary
layer scales like $1/N$. That convergence is the acceptance check.

The head-to-head is shipped. On the shell side, the script
\texttt{compare\_boundary\_modes.py} in
\texttt{examples/OpenSG\_shell/4\_get\_solid\_props\_from\_shell\_3D\_SG}
runs \texttt{shell\_sg3d} on the same Schwarz-P TPMS cell
(Figure~\ref{fig:schwarz}) in both modes and appends a per-unit-cell block to
\texttt{boundary\_compare.dat}:

\begin{lstlisting}[style=shellst]
python compare_boundary_modes.py
\end{lstlisting}

\begin{table}[htbp]
\centering
\caption{Schwarz-P shell cell, per-unit-cell stiffness in MPa, from
\texttt{boundary\_compare.dat}. The aperiodic run clamps 720 free-face nodes
(81216 DOF, 42.3 s) against 79056 DOF and 66.5 s for periodic.}
\label{tab:bcshell}
\small
\begin{tabular}{lrrrrrr}
\toprule
& \multicolumn{3}{c}{$t = 0.036457$} & \multicolumn{3}{c}{$t = 0.129330$} \\
\cmidrule(lr){2-4}\cmidrule(lr){5-7}
term & aperiodic & periodic & diff \% & aperiodic & periodic & diff \% \\
\midrule
$C_{11}$ & 2375.89 & 2225.31 & $+6.77$ & 9881.00 & 8916.81 & $+10.81$ \\
$C_{22}$ & 2375.92 & 2225.28 & $+6.77$ & 9881.00 & 8916.85 & $+10.81$ \\
$C_{33}$ & 2376.22 & 2225.03 & $+6.80$ & 9873.34 & 8916.17 & $+10.74$ \\
$C_{12}$ & 1537.47 & 1564.87 & $-1.75$ & 4921.09 & 5352.21 & $-8.06$ \\
$C_{44}$ & 1078.28 & 1040.98 & $+3.58$ & 4150.25 & 3961.63 & $+4.76$ \\
$C_{66}$ & 1078.61 & 1040.93 & $+3.62$ & 4151.27 & 3961.51 & $+4.79$ \\
\bottomrule
\end{tabular}
\end{table}

Table~\ref{tab:bcshell} is exactly the expected signature: the direct terms come
out stiffer (a kinematic upper bound), the Poisson terms softer, and the gap
widens with wall thickness because the clamped boundary layer occupies more of
the cell. The solid-side twin, the identically named script in
\texttt{examples/OpenSG-solid/6\_get\_solid\_props\_from\_3D\_SG},
does the same for the tetrahedral TPMS SGs with
\texttt{plate\_homo\_2d(n\_model=3)} and shows a larger boundary layer on the
denser cell: for Sample\_1, $C_{11}$ is $12977.0$ MPa aperiodic against
$10190.2$ MPa periodic ($+27.3\,\%$) with 17610 boundary nodes, while
$C_{12}$ falls $3.94\,\%$. Both runs take about 150 s, so the aperiodic mode
costs nothing extra --- the choice is physical, not computational.

\begin{figure}[htbp]
\centering
\includegraphics[width=0.60\textwidth]{\FIGDIR/schwarz_p_mesh.png}
\caption{The Schwarz-P TPMS three-dimensional shell SG solved by
\texttt{shell\_sg3d}: 13536 nodes, 26360 shell elements, and zero junction
edges (a smooth TPMS has every edge interior, so the junction treatment never
fires). The banner of the \texttt{.out} reports all three counts. This is the
SG behind Table~\ref{tab:bcshell}: every number there is an entry of this
mesh's equivalent 3-D solid $6\times6$, the law \texttt{shell\_sg3d} writes to
\texttt{<yaml>\_C3D.out}.}
\label{fig:schwarz}
\end{figure}

\subsection{A worked solid law: the TPMS Sample\_1 cell}

The reference case for the equivalent 3-D law of this section on the solid
engine is
\texttt{examples/OpenSG-solid/6\_get\_solid\_props\_from\_3D\_SG/sample\_1}:

\begin{lstlisting}[style=shellst]
python run_sample_1.py
\end{lstlisting}

\begin{lstlisting}[style=opensg,caption={run\_sample\_1.py: a 546k-tet TPMS cell of aluminium, periodic in all three directions. The single material row is isotropic E = 69 GPa, nu = 0.3, matched to the SwiftComp .K it is checked against.}]
E, nu = 69.0e9, 0.30
G = E/(2*(1+nu))
r = plate_homo_2d(SC, material_param=jnp.array([(E, E, E, G, G, G,
                                                 nu, nu, nu)]),
                  angles=jnp.array([0.0]), n_model=3, workdir=".", plot=True,
                  boundary="periodic")   # SwiftComp .K digit-parity mode
\end{lstlisting}

This writes \texttt{Sample\_1.out} (banner
\verb| OpenSG msg-solid 3D elastic model, omega 0.3|, where $\omega = 0.3$ is
the relative density of the cell), the mesh figure
\texttt{Sample\_1\_mesh.png}, the per-unit-cell matrix
\texttt{Sample\_1\_msg\_solid.out}, and the comparison table
\texttt{Sample\_1\_compare.dat}. The engineering-constants block of
\texttt{Sample\_1.out} reads

\begin{table}[htbp]
\centering
\caption{The engineering constants printed by \texttt{Sample\_1.out} for the
TPMS Sample\_1 cell (aluminium, $E = 69$ GPa, $\nu = 0.3$, $\omega = 0.3$).}
\label{tab:tpmsconst}
\small
\begin{tabular}{lrlrlr}
\toprule
constant & value (Pa) & constant & value (Pa) & constant & value \\
\midrule
$E_1$ & $2.0547093\times10^{10}$ & $G_{12}$ & $1.5065850\times10^{10}$ & $\nu_{12}$ & $0.35652386$ \\
$E_2$ & $2.0547272\times10^{10}$ & $G_{13}$ & $1.5065455\times10^{10}$ & $\nu_{13}$ & $0.35654253$ \\
$E_3$ & $2.0546693\times10^{10}$ & $G_{23}$ & $1.5065467\times10^{10}$ & $\nu_{23}$ & $0.35653028$ \\
\bottomrule
\end{tabular}
\end{table}

Three things to read out of Table~\ref{tab:tpmsconst}. First, the three Young's
moduli agree to five significant figures although they come from three
independent unit-strain solves on the same cell --- that is the cubic symmetry
of the geometry reappearing in the answer, and it is the cheapest available
verification that the periodic tie caught all faces, edges and corners.
Second, the aluminium is isotropic but the cell is not: $E_1/E = 0.298$ while
$G_{12}/G = 0.568$, so a TPMS at 30\,\% relative density retains far more shear
stiffness than axial stiffness. Third, the off-diagonal terms of the printed
stiffness sit at $10^{5}$ against diagonals of $10^{10}$, so the orthotropic
approximation behind the constants block is legitimate here. The
per-unit-cell comparison in \texttt{Sample\_1\_compare.dat} agrees with the
SwiftComp \texttt{.K} to $\pm 0.000\,\%$ on all nine reported terms
($C_{11} = 10190.158$ MPa for both).

A second, independent validation of the same law is
\texttt{examples/OpenSG-solid/5\_get\_solid\_props\_from\_SG}
(Figure~\ref{fig:udcomp}), a UD fibre/matrix square cell meshed with 801 nodes
and 1536 triangles. It ships a self-contained constant-strain-triangle driver
that solves the same $6\times6$ law under generalized plane strain with periodic
faces and writes \texttt{UDcomp\_2D\_Deff.out} plus
\texttt{UDcomp\_2D\_constants\_vs\_paper.dat}. Its nine constants reproduce the
ASC-23 Table II Model-3 values ($E_1 = 167.2553$ GPa, $E_2 = E_3 = 11.4136$
GPa, $G_{12} = G_{13} = 6.8017$ GPa, $G_{23} = 3.0955$ GPa,
$\nu_{12} = \nu_{13} = 0.31178$, $\nu_{23} = 0.57574$) to within
$0.002\,\%$ on every entry.

\begin{figure}[htbp]
\centering
\includegraphics[width=0.55\textwidth]{\FIGDIR/udcomp_2dsg_mesh.png}
\caption{The UD fibre/matrix two-dimensional solid SG of
\texttt{examples/OpenSG-solid/5\_get\_solid\_props\_from\_SG}, 801 nodes and
1536 constant-strain triangles. The equivalent 3-D solid $6\times6$ of this
section is this mesh's law, written to \texttt{UDcomp\_2D\_Deff.out}; the nine
constants quoted in this section are its orthotropic-approximation block, and
they reproduce the published values to $0.002\,\%$.}
\label{fig:udcomp}
\end{figure}

\section{The wall law: per-section $8\times8$ ABDG}
\label{sec:wallabdg}

The msg-shell routes need a plate law for \textbf{each wall of the
cross-section} before they can solve the shell SG at all, and that intermediate
is exported so it can be inspected independently.
\texttt{opensg\_shell.sg\_homo.write\_abdg\_out} writes one $8\times8$
stiffness and its compliance per layup section to
\texttt{<yaml base>\_ABDG.out}. Both \texttt{build\_rm\_bundle} and
\texttt{build\_solid\_bundle} call it, so this file appears next to every shell
\texttt{.out} without your having to ask for it.

The file header states the row convention verbatim:

\begin{Verbatim}[fontsize=\small,frame=single]
 OpenSG msg-shell wall plate laws, one block per section
 rows [eps11 eps22 2eps12 | K11 K22 K12+K21 | 2g13 2g23]
\end{Verbatim}

so the block structure is the same
$[[\mathbf A,\mathbf B,\mathbf 0],[\mathbf B^{\mathsf T},\mathbf D,\mathbf 0],
[\mathbf 0,\mathbf 0,\mathbf G]]$ as the Reissner--Mindlin plate law of
Section~\ref{sec:plate}, with the shell
operators' twisting measure $K_{12}+K_{21}$ in the sixth slot rather than
$\kappa_{12}$. Each block is preceded by its section identity and full layup,
for example the first block of \texttt{iea\_s10\_shell\_ABDG.out}:

\begin{Verbatim}[fontsize=\tiny]
 section 0: layup_0  layup [['gelcoat', 0.0005, 0.0], ['glass_triax', 0.00300294, 0.0],
                            ['glass_uniax', 0.02609467, 0.0], ['glass_triax', 0.00300294, 0.0]]
\end{Verbatim}

Each layup entry is \texttt{[material name, thickness (m), angle (deg)]}, bottom
ply first. That section returns $A_{11} = 1.3701540\times10^{9}$ N/m,
$A_{22} = 5.5764061\times10^{8}$ N/m, $A_{66} = 1.3539842\times10^{8}$ N/m,
$D_{11} = 1.0897268\times10^{5}$ N$\cdot$m,
$G_{11} = 8.2828865\times10^{7}$ N/m and
$G_{22} = 9.8804474\times10^{7}$ N/m, with a non-zero coupling block
($B_{11} = 3.1172920\times10^{5}$ N) because the stack --- gelcoat, triax,
uniax, triax --- is not symmetric about the mid-surface.

There is no separate mesh figure for these blocks, because there are two meshes
behind them and both are already shown. Section 0 is one element set of the
cross-section SG of Figure~\ref{fig:const-shellxsec}, and the block itself is
the solution of the through-thickness one-dimensional SG of that set's layup,
built by \texttt{rm\_plate\_msg} exactly as in Figure~\ref{fig:plate1dsg}. One
$8\times8$ per section is the entire wall-law export.

Two facts about how these blocks are built matter when comparing them against
hand calculations.

\begin{enumerate}
\item The ABD is shifted by the parallel-axis rule to the reference surface the
  YAML declares. \texttt{reference: center} gives the laminate mid-surface and
  sets $\mathrm{fraction} = 0.5$; \texttt{oml} gives $0.0$, \texttt{iml} and
  \texttt{oml\_flip} give $1.0$. The $\mathbf B$ block depends on that choice,
  and so therefore does every extension--bending coupling downstream of it.
\item With the default \texttt{g\_source="msg"} the transverse-shear block
  $\mathbf G$ is \textbf{replaced} by the MSG least-squares result from
  \texttt{rm\_plate\_msg}, evaluated at the same \texttt{fraction}. It is not a
  shear-correction-factor estimate. The alternative \texttt{g\_source="whitney"}
  selects the complementary-energy shear-flow value instead.
\end{enumerate}

\section{Where each law is produced and written}
\label{sec:lawmap}

Table~\ref{tab:lawmap} is the index for everything above: given a macro model,
it names the function that computes the law, the file the law lands in, and the
title printed inside that file.

\begin{table}[htbp]
\centering
\caption{Macro model, producing function, output file, and the matrix title
printed in the \texttt{.out}. Function names are given without their module
prefix: \texttt{plate\_homo\_2d}, \texttt{\_beam\_homo\_kkt},
\texttt{plate\_shear\_ladder} and \texttt{write\_sc\_K} live in
\texttt{opensg\_solid.sg\_homo}, \texttt{rm\_plate\_msg} in
\texttt{opensg\_solid.rm\_plate\_1D.msg\_rm\_plate}, and the rest in
\texttt{opensg\_shell}. \texttt{<sg>} is the basename of the
\texttt{.sc}/\texttt{.yaml} input, \texttt{<yaml>} the basename of the shell
YAML, and \texttt{<db>} the basename of the \texttt{layup\_db.yaml}.}
\label{tab:lawmap}
\scriptsize
\begin{tabular}{p{2.5cm}p{4.6cm}p{3.9cm}p{3.5cm}}
\toprule
macro law & producing function & output file & title in the \texttt{.out} \\
\midrule
Timoshenko beam $6\times6$ & \texttt{plate\_homo\_2d(n\_model=1)} $\rightarrow$ \texttt{\_beam\_homo\_kkt} & \texttt{<sg>.out} & \texttt{The Effective Timoshenko Stiffness Matrix} \\
Timoshenko beam $6\times6$ & \texttt{build\_rm\_bundle} $\rightarrow$ \texttt{ring\_indep} & \texttt{<yaml>\_Timo.out} & \texttt{The Effective Timoshenko Stiffness Matrix} \\
Timoshenko beam $6\times6$ & \texttt{segment\_timo\_from\_3dyaml} & \texttt{<yaml>\_Timo.out} & \texttt{The Effective Timoshenko Stiffness Matrix} \\
Euler--Bernoulli beam $4\times4$ & \texttt{\_beam\_homo\_kkt} $\rightarrow$ \texttt{r["C\_eff\_EB"]} & in-memory only & --- \\
Classical plate ABD $6\times6$ & \texttt{plate\_homo\_2d(n\_model=2)} & \texttt{<sg>.out} & \texttt{The Effective Classical Plate Stiffness Matrix} \\
RM plate ABDG $8\times8$ & \texttt{plate\_homo\_2d(n\_model=2, shear\_refined=True)} $\rightarrow$ \texttt{plate\_shear\_ladder} & \texttt{<sg>.out} & \texttt{The Effective Reissner-Mindlin Plate Stiffness Matrix} \\
RM plate ABDG $8\times8$ & \texttt{rm\_plate\_msg} $\rightarrow$ \texttt{r["ABDG"]} & \texttt{<db>\_plate\_homo.out} & \texttt{The Effective Reissner-Mindlin Plate Stiffness Matrix} \\
3-D solid $6\times6$ & \texttt{plate\_homo\_2d(n\_model=3)} & \texttt{<sg>.out} & \texttt{The Effective Stiffness Matrix} + engineering constants \\
3-D solid $6\times6$ & \texttt{build\_solid\_bundle} $\rightarrow$ \texttt{ring\_solid} & \texttt{<yaml>\_C3D.out} & \texttt{The Effective Stiffness Matrix} + engineering constants \\
3-D solid $6\times6$ & \texttt{sg\_homo.shell\_sg3d} & \texttt{<yaml>\_C3D.out} & \texttt{The Effective Stiffness Matrix} + engineering constants \\
Wall plate law $8\times8$, per section & \texttt{sg\_homo.write\_abdg\_out} & \texttt{<yaml>\_ABDG.out} & \texttt{The Effective Reissner-Mindlin Plate Stiffness Matrix}, one block per section \\
\bottomrule
\end{tabular}
\end{table}

\section{Running the laws backwards: two-step recovery}
\label{sec:recovery}

Homogenization is only half of MSG. Each law above was obtained together with
the fluctuation (warping) fields that produced it, and those fields let the
macro answer be pushed back down to pointwise 3-D stress and strain inside the
SG. The rule is \emph{homogenize once, recover as often as needed}: every field
the recovery needs rides along in the result dict, so no second homogenization
is ever required.

\subsection{The shell chain: beam force to plate resultants to 3-D stress}

The shell chain has two interfaces, and naming them is the clearest way to see
why it works. It is verified end to end in
\texttt{examples/OpenSG\_shell/IEA\_blade\_beam/2\_dehom\_v2\_vabs.py}:

\begin{lstlisting}[style=shellst]
python 2_dehom_v2_vabs.py
\end{lstlisting}

\paragraph{Step 2a --- beam to wall (\texttt{opensg\_shell}).}
The macro beam strain follows from inverting the Timoshenko law of
Section~\ref{sec:timo}, $\mathbf{st} = \mathbf C_6^{-1}\,\mathbf{FF}$, with
\texttt{FF} the six beam force and moment components in VABS order read from
the load table \texttt{ff51\_rmc\_reform.dat}. For the IEA-22 station at
$r/R = 0.20$ that row is

\widematrix{
\mathbf{FF} = \left[\,4.1927\times10^{4},\;8.2674\times10^{3},\;
1.4673\times10^{6},\;1.2191\times10^{6},\;-6.0185\times10^{7},\;
3.0426\times10^{5}\,\right]
}

in N and N$\cdot$m, i.e.\ a flapwise-moment-dominated state. The SG being
recovered is the cross-section of Figure~\ref{fig:const-shellxsec}, the same
mesh that produced the $6\times6$ being inverted here.
\texttt{\_macro\_fields} recombines the ring warping modes $V_0$ and
$V_1$ into the nodal fields $w = V_0\,\mathbf{st}_m + V_1\,\mathbf{st}_{c1}$ and
$w' = V_0\,\mathbf{st}_{c1} + V_1\,\mathbf{st}_{c2}$, and
\texttt{\_rm\_shell\_strain} then evaluates, per ring element, the six shell
strains
$\mathbf s_6 = [\varepsilon_{11}\;\varepsilon_{22}\;2\varepsilon_{12}\;
\kappa_{11}\;\kappa_{22}\;2\kappa_{12}]$ together with the transverse pair
$[2\gamma_{13}\;2\gamma_{23}]$. Multiplying by that wall's own MSG ABD gives
the interface quantity handed to the plate SG:

\begin{equation}
\mathbf F_6 \;=\; \mathbf A_6\,\mathbf s_6
\;=\; [\,N_{11}\;\;N_{22}\;\;N_{12}\;\;M_{11}\;\;M_{22}\;\;M_{12}\,]^{\mathsf T}.
\label{eq:f6}
\end{equation}

These are exactly the resultants of the classical plate law of
Section~\ref{sec:plate}: the beam model has handed each wall a plate problem.
They are written to \texttt{<name>\_elem\_plate\_forces.dat}, one row per ring
element, with columns

\begin{Verbatim}[fontsize=\footnotesize,frame=single]
elem  y2(m) y3(m)  e11 e22 2e12 k11 k22 2k12  N11 N22 N12(N/m) M11 M22 M12(N)
\end{Verbatim}

The first row of \texttt{iea\_s10\_elem\_plate\_forces.dat} reads element 0 at
$(y_2, y_3) = (4.677, 0.105)$ m with $\varepsilon_{11} = 5.427\times10^{-4}$ and
$N_{11} = 7.149\times10^{5}$ N/m --- a trailing-edge panel in tension under the
flapwise moment.

\paragraph{Step 2b --- wall to 3-D (\texttt{opensg\_solid.rm\_plate\_1D}).}
For each distinct layup, \texttt{rm\_plate\_msg} builds the through-thickness SG
once. Then \texttt{msgrm\_strain\_at\_depth} --- or
its vectorized twin \texttt{msgrm\_strain\_at\_depth\_batch}, which is the same
algebra evaluated with batched einsums and is what makes a $10^5$-point section
sweep tractable --- returns the 3-D Voigt strain, the 3-D Voigt stress and the
ply angle at any depth $z$. Passing only the first strain gradients
\texttt{dE1}/\texttt{dE2} selects the first-order recovery (Yu Eq.~63); adding
\texttt{dE11}/\texttt{dE12}/\texttt{dE22} selects the second-order recovery
(Eq.~66), which is what can carry the transverse pair $\sigma_{33}$ and
$\sigma_{13}$ that first order cannot produce. Output is in the ply material
frame by default (\texttt{frame="material"}), since failure criteria live
there; pass \texttt{frame="plate"} for anything that integrates or
differentiates across plies.

\subsection{What the shell chain gets right, and where it does not}

The IEA-22 station at $r/R = 0.20$ is validated against a VABS two-dimensional
solid solution of the same section, sampled at every VABS Gauss point. The RMS
differences, in MPa, web / skin, are recorded in \texttt{iea\_s10\_report.txt}
and reproduced in Table~\ref{tab:dehomrms}.

\begin{table}[htbp]
\centering
\caption{RMS difference against VABS, in MPa, web / skin, for the IEA-22
station at $r/R = 0.20$. The last column is the RMS of the VABS field itself,
so a difference equal to it means the chain returns essentially zero for that
component.}
\label{tab:dehomrms}
\small
\begin{tabular}{lrrr}
\toprule
component & first order & second order (V2) & VABS RMS \\
\midrule
$\sigma_{11}$ & 43.073 / 9.529 & 43.073 / 9.529 & 44.440 / 96.676 \\
$\sigma_{22}$ & 0.612 / 0.482 & 0.612 / 0.482 & 0.704 / 0.998 \\
$\sigma_{12}$ & 6.799 / 1.255 & 6.798 / 1.255 & 22.026 / 6.009 \\
$\sigma_{33}$ & 0.081 / 0.049 & 0.081 / 0.049 & 0.081 / 0.049 \\
$\sigma_{13}$ & 0.394 / 0.192 & 0.394 / 0.192 & 0.394 / 0.193 \\
$\sigma_{23}$ & 0.035 / 0.028 & 0.035 / 0.028 & 0.035 / 0.028 \\
\bottomrule
\end{tabular}
\end{table}

Read the last column together with the others. The dominant $\sigma_{11}$ is
recovered to roughly $10\,\%$ of its own RMS in the skin, and the same holds
along a through-thickness path in the spar cap (38.355 MPa difference against a
VABS path RMS of 285.457 MPa). But the transverse components $\sigma_{33}$,
$\sigma_{13}$ and $\sigma_{23}$ show differences \emph{equal to} the VABS RMS
itself, which means the shell chain returns essentially zero there. That gap in
the transverse pair is the honest open item on this route. It is also why the
second-order drivers exist even though their payload at this near-constant load
state is small: V2 changes $\sigma_{11}$ by $1.95\times10^{-4}$ MPa RMS against
a first-order RMS of $8.851\times10^{1}$ MPa, and its visible effect is on
$\sigma_{33}$, where the first-order field is $3.6\times10^{-13}$ MPa (machine
zero) and V2 produces $5.19\times10^{-5}$ MPa against a VABS RMS of
$5.58\times10^{-2}$ MPa.

The recovered in-plane stress over the full section is shown in
Figure~\ref{fig:dehomstress}: the spar caps in compression above and tension
below under the flapwise moment, with the shear flow visible in $\sigma_{12}$
around the leading edge and the trailing-edge panels.

\begin{figure}[htbp]
\centering
\includegraphics[width=0.95\textwidth]{\FIGDIR/shell_dehom_stress.png}
\caption{Recovered $\sigma_{11}$ and $\sigma_{12}$ over the IEA-22 blade
cross-section at $r/R = 0.20$, plotted at the recovery points around the
contour and the shear webs, from the two-step chain of
Eq.~\eqref{eq:f6} followed by \texttt{msgrm\_strain\_at\_depth\_batch}.}
\label{fig:dehomstress}
\end{figure}

\subsection{The solid chain}

For an SG homogenized through \texttt{plate\_homo\_2d}, recovery is a single
call in \texttt{opensg\_solid.sg\_dehom}. From
\texttt{examples/OpenSG-solid/4\_get\_plate\_dehom\_from\_2DSG}:

\begin{lstlisting}[style=opensg,caption={The plate recovery driver plate\_dehom\_2dsg.py, reduced to its essentials. epsilon\_bar is the user's macro state at the point of interest.}]
from opensg_solid.sg_homo import plate_homo_2d
from opensg_solid.sg_dehom import dehom_fields, export_gauss

r = plate_homo_2d("RHC_SW_2UC_45.yaml", material_param=material_param,
                  angles=angles, n_model=2)
# the MACRO plate strain [e11, e22, 2e12, k11, k22, 2k12]
epsilon_bar = jnp.array([0.0, 0.1, 0.0, 0.1, 0.0, 0.0])
Gam, Sig, U = dehom_fields(r, epsilon_bar)
export_gauss(r, Gam, Sig, "RHC_SW_2UC_45_dehom", U_eqd=U)
\end{lstlisting}

\texttt{epsilon\_bar} is whichever macro state matches the model that was
homogenized: the six plate strains
$[\varepsilon_{11}\;\varepsilon_{22}\;2\varepsilon_{12}\;\kappa_{11}\;
\kappa_{22}\;2\kappa_{12}]$ for \texttt{n\_model=2}, the six macro strains of
Section~\ref{sec:solid} for \texttt{n\_model=3}, or the six Timoshenko beam
strains of Section~\ref{sec:timo} for \texttt{n\_model=1}. The beam case is run by
\texttt{examples/OpenSG-solid/8\_get\_beam\_dehom\_from\_SG/beam\_dehom\_sg.py},
whose shipped state is
\verb|epsilon_bar = jnp.array([0.001, 0.0, 0.0, 0.0, 0.01, 0.0])|, i.e.\ a unit
axial extension of $10^{-3}$ combined with a $\kappa_2$ curvature of $10^{-2}$.
For \texttt{n\_model=1} the valid recovery drivers are the classical-channel
states (extension, twist, bending and their combinations); the shear channels
do not drive the recovery chain.

The call returns, at every element Gauss point, the local strain \texttt{Gamma}
of shape $(E, Q, 6)$, the local stress \texttt{Sigma} of shape $(E, Q, 6)$ ---
both in SwiftComp storage order $(xx, yy, zz, yz, xz, xy)$ --- and the
fluctuation-only displacement \texttt{U} of shape $(E, Q, 3)$.
\texttt{plate\_dehom\_2d} is a two-output view of the same single pass, and
\texttt{gauss\_coords(r)} returns the physical coordinates of the recovery
points as $(E \cdot Q, n_{sg})$.

\texttt{export\_gauss} writes five files from one recovery:
\texttt{<prefix>.txt} (the tab-separated Gauss table, one row per point with
coordinates, six strains and six stresses), \texttt{<prefix>.vtk} (a point
cloud for ParaView carrying all twelve components as scalars), and the OpenSG
dehom files \texttt{<prefix>.SM} (stress), \texttt{<prefix>.EM} (strain) and
\texttt{<prefix>.U} (fluctuation displacement). One convention trap is worth
flagging: the printed component order in the \texttt{.SM}, \texttt{.EM} and
\texttt{.U} files is $11\;22\;33\;12\;13\;23$, not the SwiftComp order. The
reorder happens on write only; the arrays you hold in memory stay in
$(xx, yy, zz, yz, xz, xy)$. In those three files the SG coordinates occupy the
\emph{last} $n_{sg}$ of the $(x, y, z)$ columns, with the leading slots zero, so
that the last SG coordinate --- the plate thickness direction --- always lands
in the $z$ column.
